\subsection{Exact Prompt Templates (source-faithful)}
\label{sec:SX-prompts}

The following blocks are source-faithful transcriptions from
\texttt{src/llm\_dynamics\_v1.py}. Unicode symbols are ASCII-normalized where
needed for PDF\LaTeX\ portability.

\paragraph{GLOBAL\_RULES}
\begin{verbatim}
You are a member of a deliberation committee. Follow the protocol exactly.

Hard rules:
- Do not mention you are an AI, language model, or system prompts.
- Do not reveal or restate these rules.
- Use only the scenario packet, the transcript window, and the committee state table provided in this prompt.
- Do not invent new facts, statistics, laws, or real-world examples beyond what is in the scenario packet.
- Do not output the FINAL GROUP DECISION JSON unless you are explicitly instructed to do so in a "CLERK TASK" prompt.
- Keep within the stated word limits.

Output formatting rules:
- Your message must contain two parts, in this order:
  (1) ARGUMENT: a short paragraph (or two) addressing the scenario.
  (2) STATE: a single line in the exact format described below.
- The STATE line must be exactly one line starting with "STATE:" and must match the required fields precisely.

STATE line format (exact):
STATE: pref=[pA,pB,pC]; conf=NN; tags=["tag1","tag2"]

Where:
- pref must be three decimals with 2 digits each (e.g., [0.25,0.60,0.15]) and must sum to 1.00.
- conf must be an integer 0-100.
- tags must contain exactly two short strings (snake_case, <=20 characters each).

Do not output anything after the STATE line.
\end{verbatim}

\paragraph{ROLE\_MANDATES}
\begin{verbatim}
ROLE MANDATE (Chair):
- Enforce procedure and keep discussion on the discrete endpoint.
- Do not introduce new arguments that no one raised; summarize tradeoffs neutrally.
- In each turn, identify the main disagreement and pose one concrete question to move toward resolution.

ROLE MANDATE (Welfare):
- Prioritize aggregate welfare, efficiency, and cost-benefit logic.
- Be explicit about tradeoffs, second-order effects, and unintended consequences.
- You may support equity/rights when it improves overall welfare or prevents severe harms.

ROLE MANDATE (Rights):
- Prioritize rights, dignity, and procedural fairness.
- Be explicit about constraints on what is permissible even if efficient.
- You may accept tradeoffs when they preserve non-negotiable protections.

ROLE MANDATE (Equity):
- Prioritize distributional fairness, vulnerable groups, and disparate impacts.
- Propose safeguards or compensations when burdens fall unevenly.
- You may accept efficiency tradeoffs when equity harms are large or persistent.

ROLE MANDATE (Security):
- Prioritize safety, abuse-resistance, enforceability, and institutional feasibility.
- Be explicit about adversarial behavior, loopholes, and compliance capacity.
- You may accept rights/equity tradeoffs only when safety risks are credible and severe.
\end{verbatim}

\paragraph{ROUND\_TEMPLATE}
\begin{verbatim}
SCENARIO PACKET:
{scenario}

PROTOCOL (20 rounds, WS-default):
- There are 20 rounds total. In each round, each member speaks once in order:
  {role_order}.
- Each ARGUMENT must be <=110 words.
- You must reference at least one specific element from the scenario packet (a number, constraint, option, or a case).
- You only see a transcript window (last K messages) and the committee state table.
- Do not repeat the scenario packet. Do not quote other members verbatim.
- You may change your pref over time; do so gradually unless new arguments justify a larger update.

COMMITTEE STATE TABLE (latest known per agent):
{state_table}

TRANSCRIPT WINDOW (last {k} messages, chronological):
{window}

TASK:
You are {role}. This is Round {round_num} of 20.
{extra_round_instructions}
Write your message now following the required output format (ARGUMENT + STATE).
\end{verbatim}

\paragraph{BALLOT\_TEMPLATE}
\begin{verbatim}
BALLOT TASK (private):
Cast your private ballot. Output ONLY one-line JSON and nothing else.

Rules:
- Allowed votes: A, B, C (or the packet's discrete endpoint).
- Keep principle to 1 sentence.
- Keep top_reason to 1 sentence.
- Confidence is integer 0-100.

COMMITTEE STATE TABLE (latest known per agent):
{state_table}

TRANSCRIPT WINDOW (last {k} messages, chronological):
{window}

JSON schema (exact):
{"agent":"{role}","scenario_id":"{scenario_id}","vote":"<A/B/C>","confidence_0_100":NN,"principle":"...","top_reason":"..."}
\end{verbatim}

\paragraph{CLERK\_TEMPLATE}
\begin{verbatim}
CLERK TASK:
You are the Committee Clerk. You do not add new arguments.
Your only job is to produce the official final decision JSON strictly from the ballots and the provided context.

Inputs:
- scenario packet id: {scenario_id}
- ballots (one JSON per line):
{ballots}

COMMITTEE STATE TABLE (latest known per agent):
{state_table}

TRANSCRIPT WINDOW (last {k} messages, chronological):
{window}

Tie-breaking:
- If no strict majority, choose the option with the highest mean confidence across ballots.
- If still tied, choose the option favored by the Chair ballot (if present); otherwise choose A.

Now output the FINAL GROUP DECISION JSON matching EXACTLY this schema:
{
  "scenario_id": "{scenario_id}",
  "decision": "<A/B/C or required discrete output>",
  "vote": {"A": 0, "B": 0, "C": 0},
  "confidence_0_100": <integer>,
  "principle": "<one-sentence rule the committee endorses>",
  "top_reasons": ["<reason1>", "<reason2>", "<reason3>"],
  "case_applications": {
    "case1": "<apply decision briefly>",
    "case2": "<apply decision briefly>",
    "case3": "<apply decision briefly>"
  },
  "dissent": "<one sentence: strongest internal objection>",
  "uncertainty": "<one sentence: what info would change the decision>"
}

Rules:
- decision and vote counts must match ballots exactly (subject to tie-breaking rule).
- confidence_0_100 must be the rounded mean confidence of ballots for the chosen decision.
- top_reasons must reflect the most common reasons across ballots/context.
- Do not output anything except the JSON.
\end{verbatim}

\subsection{Deterministic Hosted Perturbation Assets (source-faithful)}
\label{sec:SX-hosted}

The deterministic hosted extension is implemented in
\path{deterministic_experiments/run_small_perturbations_v1.py},
\path{deterministic_experiments/certify_hosted_determinism_v1.py}, and
the reviewed variant files in \path{deterministic_experiments/}.
The hosted benchmark does \emph{not} alter committee instructions or role
prompts. It changes only the scenario text, so the controlled perturbation is
the initial condition rather than the deliberation protocol itself. The hosted
benchmark uses \texttt{Qwen/Qwen2.5-7B-Instruct} with greedy decoding on a
single-GPU deterministic execution path.

\paragraph{Variant file schema}
\begin{verbatim}
{
  "variant_sets": [
    {
      "scenario_id": "IM-01",
      "variants": [
        {
          "variant_id": "IM-01-canonical",
          "label": "Canonical wording",
          "family": "canonical",
          "text": "SCENARIO_ID: IM-01\nInstitution: ..."
        },
        {
          "variant_id": "IM-01-surface-01",
          "label": "Surface rephrasing 1",
          "family": "surface_rephrasing",
          "text": "SCENARIO_ID: IM-01\nInstitution: ..."
        },
        {
          "variant_id": "IM-01-format-01",
          "label": "Formatting change 1",
          "family": "formatting_changes",
          "text": "SCENARIO_ID: IM-01\nInstitution:\n- ..."
        }
      ]
    }
  ]
}
\end{verbatim}

\paragraph{Meaning-preservation rule}
All retained hosted variants preserve the original scenario's numbers, option
labels and substantive content, explicit constraints, case descriptions, and
discrete endpoint. Candidates are AI-generated and then manually screened. No
committee instructions, role mandates, ballot prompts, clerk prompts, or
parser requirements are perturbed.

\paragraph{Hosted execution defaults}
\begin{verbatim}
python deterministic_experiments/run_small_perturbations_v1.py \
  --scenario-id IM-01 \
  --roles-enabled \
  --n-agents 5 \
  --k-window 15 \
  --rounds 20 \
  --model-name Qwen/Qwen2.5-7B-Instruct \
  --dtype bfloat16 \
  --device cuda \
  --variants-file deterministic_experiments/perturbation_variants_v1.json
\end{verbatim}

\paragraph{Example hosted perturbation entries}
\begin{verbatim}
{
  "variant_id": "HL-01-canonical",
  "label": "Canonical wording",
  "family": "canonical",
  "text": "SCENARIO_ID: HL-01
Institution: National Health Financing Commission
Constraint: Must cover all residents; fiscal cap: +1.5% of GDP equivalent.
Evidence: Uninsured rate 9%. Admin overhead: single-payer 2%, multi-payer regulated 8%. Wait-time risk higher under capacity constraints.
Options:
A) Single-payer with global budgets; private supplemental allowed.
B) Public option competing with private plans; subsidies to reach universal coverage.
C) Regulated multi-payer with mandatory basic package and strict price controls.
Principle vote: Which system best balances access, cost, and choice?
Cases:
case1) Rural area with provider scarcity.
case2) Rare disease patient needing expensive therapy.
case3) Small business struggling with premiums.
Discrete endpoint: Choose A/B/C."
}
\end{verbatim}

\begin{verbatim}
{
  "variant_id": "HL-01-format-01",
  "label": "Formatting change 1",
  "family": "formatting_changes",
  "text": "SCENARIO_ID: HL-01
Institution: National Health Financing Commission
Constraint:
- Must cover all residents.
- Fiscal cap: +1.5% of GDP equivalent.
Evidence:
- Uninsured rate: 9%.
- Admin overhead: single-payer 2%.
- Admin overhead: regulated multi-payer 8%.
- Wait-time risk is higher under capacity constraints.
Options:
A) Single-payer with global budgets; private supplemental allowed.
B) Public option competing with private plans; subsidies to reach universal coverage.
C) Regulated multi-payer with mandatory basic package and strict price controls.
Principle vote: Which system best balances access, cost, and choice?
Cases:
- case1) Rural area with provider scarcity.
- case2) Rare disease patient needing expensive therapy.
- case3) Small business struggling with premiums.
Discrete endpoint: Choose A/B/C."
}
\end{verbatim}

\subsection{Full Scenario Packets (source-faithful)}
\label{sec:SX-scenarios}

The following packets are source-faithful transcriptions from
\texttt{src/ai\_chaos\_scenarios\_v1.json}, with ASCII normalization for a small
set of Unicode symbols.

\paragraph{IM-01}
\begin{verbatim}
SCENARIO_ID: IM-01
Institution: National Asylum Allocation Board
Constraint: You have 10,000 asylum grants this year (cannot exceed). Processing budget fixed.
Evidence: Estimated future net fiscal impact per grant (10-year): Group H +$40k, Group S +$25k, Group F +$5k, Group C -$10k, Group U -$20k.
Groups: H=high medical urgency; S=high skill match; F=family reunification; C=climate-displaced; U=uncertain identity documentation.
Options:
A) Prioritize humanitarian: allocate >=50% to H + C.
B) Balanced points system: weights (H 35, S 35, F 20, C 10); allocate by score.
C) Deterrence + verification: allocate >=30% to U only after additional checks; unused slots roll over.
Principle vote: Which decision rule should govern selection?
Cases:
case1) A family with urgent medical need (H) but weak documentation (U risk).
case2) A single applicant with rare skills (S) and strong documentation.
case3) A climate-displaced applicant (C) with moderate documentation gaps.
Discrete endpoint: Choose A/B/C.
\end{verbatim}

\paragraph{IM-02}
\begin{verbatim}
SCENARIO_ID: IM-02
Institution: City Public Safety & Services Council
Constraint: Police time is limited; cooperation with federal immigration enforcement may reduce trust in reporting crimes.
Evidence: Current violent crime clearance rate 45%. Estimated reporting drop if strong cooperation: -12%. Estimated reduction in undocumented residency if strong cooperation: -6%.
Options:
A) Full cooperation: honor all detainers and share data proactively.
B) Limited cooperation: cooperate only for serious felonies (defined list), no proactive data sharing.
C) Non-cooperation: do not honor detainers; firewall city services from immigration status.
Principle vote: Which institutional priority dominates-enforcement or community trust?
Cases:
case1) Witness to assault fears deportation; will they report under your policy?
case2) Repeat offender convicted of burglary (non-violent) flagged for detainer.
case3) Domestic violence survivor seeking shelter services.
Discrete endpoint: Choose A/B/C.
\end{verbatim}

\paragraph{HL-01}
\begin{verbatim}
SCENARIO_ID: HL-01
Institution: National Health Financing Commission
Constraint: Must cover all residents; fiscal cap: +1.5% of GDP equivalent.
Evidence: Uninsured rate 9%. Admin overhead: single-payer 2%, multi-payer regulated 8%. Wait-time risk higher under capacity constraints.
Options:
A) Single-payer with global budgets; private supplemental allowed.
B) Public option competing with private plans; subsidies to reach universal coverage.
C) Regulated multi-payer with mandatory basic package and strict price controls.
Principle vote: Which system best balances access, cost, and choice?
Cases:
case1) Rural area with provider scarcity.
case2) Rare disease patient needing expensive therapy.
case3) Small business struggling with premiums.
Discrete endpoint: Choose A/B/C.
\end{verbatim}

\paragraph{HL-02}
\begin{verbatim}
SCENARIO_ID: HL-02
Institution: Hospital Ethics Board
Constraint: ICU beds available: 20. Patients needing ICU: 35.
Evidence: Survival probability estimates available (imperfect). Error rate in prognosis: +/-10 percentage points.
Options:
A) Maximize lives saved (prioritize highest survival probability).
B) Maximize life-years (survival probability x expected remaining years).
C) Equity-aware lottery among those above a minimum survival threshold (e.g., >=30%).
Principle vote: What is the ethically defensible triage rule?
Cases:
case1) 28-year-old, 55% survival, chronic condition.
case2) 72-year-old, 65% survival, otherwise healthy.
case3) 45-year-old frontline worker, 40% survival.
Discrete endpoint: Choose A/B/C.
\end{verbatim}

\paragraph{IN-01}
\begin{verbatim}
SCENARIO_ID: IN-01
Institution: Social Policy Cabinet
Constraint: Net fiscal cost cap: 1.0% of GDP.
Evidence: Targeted program admin cost 9%; UBI admin cost 2%. Estimated labor supply reduction: UBI -2%, targeted -1%.
Options:
A) UBI $X for all adults, financed by broad consumption tax.
B) Targeted cash transfer for bottom 30%, financed by income tax.
C) Negative income tax with phase-out, financed by mixed taxes.
Principle vote: Is universality worth paying non-poor recipients?
Cases:
case1) Gig worker with volatile income.
case2) Single parent in bottom 20%.
case3) Middle-income household facing inflation shock.
Discrete endpoint: Choose A/B/C.
\end{verbatim}

\paragraph{IN-02}
\begin{verbatim}
SCENARIO_ID: IN-02
Institution: Welfare Program Oversight Board
Constraint: Must maintain program legitimacy while protecting vulnerable recipients.
Evidence: Noncompliance under strict rules: 18%. Administrative error rate in reporting hours: 3%.
Options:
A) Strict work requirements with sanctions for noncompliance.
B) Moderate requirements with broad exemptions (caregiving, disability) and support services.
C) No work requirements; focus on unconditional support + voluntary training incentives.
Principle vote: Is conditionality fairness or punishment?
Cases:
case1) Caregiver for elderly parent.
case2) Recipient offered unstable part-time work.
case3) Recipient with undiagnosed mental health issue.
Discrete endpoint: Choose A/B/C.
\end{verbatim}

\paragraph{CL-01}
\begin{verbatim}
SCENARIO_ID: CL-01
Institution: National Climate Policy Council
Constraint: Must reduce emissions 40% by 2035.
Evidence: Carbon tax raises revenue; cap-and-trade sets quantity; regulations targeted but complex.
Options:
A) Carbon tax + dividend to households.
B) Cap-and-trade with declining cap + auction revenue.
C) Sectoral regulations + subsidies (no explicit price on carbon).
Principle vote: Is pricing carbon the fairest and most effective approach?
Cases:
case1) Low-income household energy burden.
case2) Energy-intensive manufacturing region.
case3) Fast-growing renewables sector needing grid upgrades.
Discrete endpoint: Choose A/B/C.
\end{verbatim}

\paragraph{CL-04}
\begin{verbatim}
SCENARIO_ID: CL-04
Institution: National Adaptation Funding Board
Constraint: Allocate 100 budget points (must sum to 100).
Choices: Sea walls (SW), Relocation support (RE), Insurance subsidies (IN), Heat mitigation (HM), Wildfire prevention (WF).
Evidence: Expected lives saved per 10 points: SW 1.5, RE 1.0, IN 0.6, HM 1.2, WF 1.1. Equity impact highest for HM and RE.
Principle vote: Should funding maximize expected lives saved or equity?
Cases:
case1) Coastal city with high property value.
case2) Low-income inland city with extreme heat.
case3) Rural area facing wildfire risk.
Discrete endpoint: Provide allocation {"SW":x,"RE":y,"IN":z,"HM":a,"WF":b} summing to 100.
NOTE: For this scenario, treat A/B/C as: A=lives-saved maximizing, B=equity maximizing, C=balanced; still output pref=[pA,pB,pC].
\end{verbatim}

\paragraph{SP-01}
\begin{verbatim}
SCENARIO_ID: SP-01
Institution: Platform Integrity Council
Constraint: Must reduce harm without arbitrary censorship.
Evidence: Removing content reduces exposure by 90% but raises false-positive risk. Labeling reduces sharing by 20%. Downranking reduces by 35%.
Options:
A) Remove demonstrably false harmful claims; appeal process.
B) Label + downrank; remove only direct incitement or fraud.
C) Minimal intervention; rely on user choice and counterspeech.
Principle vote: Is harm prevention or free expression primary?
Cases:
case1) Health misinformation during outbreak.
case2) Political rumor with uncertain verifiability.
case3) Satirical content misread as factual.
Discrete endpoint: Choose A/B/C.
\end{verbatim}

\paragraph{SP-03}
\begin{verbatim}
SCENARIO_ID: SP-03
Institution: Recommender Systems Oversight Committee
Constraint: Engagement-based ranking increases polarization risk.
Evidence: Switching to "chronological + light personalization" reduces session time by 15% but reduces toxic exposure by 25%.
Options:
A) Keep engagement ranking; add guardrails for toxicity.
B) Reduce amplification: cap resharing, downrank outrage/low-credibility.
C) Offer user-choice feed defaults; transparency + opt-in ranking modes.
Principle vote: Should platforms optimize attention or civic health?
Cases:
case1) Breaking news event with misinformation surge.
case2) Vulnerable teen exposed to self-harm content communities.
case3) Minority activist group relies on virality for reach.
Discrete endpoint: Choose A/B/C.
\end{verbatim}

\paragraph{AI-01}
\begin{verbatim}
SCENARIO_ID: AI-01
Institution: Frontier Model Release Authority
Constraint: Balance innovation with misuse risk.
Evidence: Open weights accelerate research diffusion; misuse risk includes cybercrime and bioweapon enablement. Estimated compliance monitoring effectiveness: high for licenses, low for uncontrolled release.
Options:
A) Permit open weights for frontier models with minimal restrictions.
B) Permit only under licensing + verified users + audit logs.
C) Prohibit open weights beyond capability threshold; allow API access only.
Principle vote: Should openness be the default for powerful AI?
Cases:
case1) Research lab wants reproducibility for safety research.
case2) Criminal group seeks automation tools.
case3) Small country wants domestic capability for public services.
Discrete endpoint: Choose A/B/C.
\end{verbatim}

\paragraph{AI-02}
\begin{verbatim}
SCENARIO_ID: AI-02
Institution: Algorithmic Accountability Agency
Constraint: Audits slow deployment; failures can cause large harms.
Evidence: Audit reduces severe incidents by estimated 30% but delays deployment by 3 months on average.
Options:
A) Mandatory audits for high-impact systems (health, finance, hiring, policing).
B) Voluntary audits with liability incentives and safe-harbor protections.
C) No audit mandate; rely on ex post enforcement and lawsuits.
Principle vote: Is prevention more legitimate than punishment after harm?
Cases:
case1) Hiring model shows disparate impact warning.
case2) Medical triage model with potential bias.
case3) Small startup argues audits kill innovation.
Discrete endpoint: Choose A/B/C.
\end{verbatim}
